\textbf{Phần 1. Trắc nghiệm lựa chọn}
\setcounterpageref{ex}{0}

\begin{ex} %[5X11N65]
Tung một đồng xu $20$ lần, thấy có $11$ lần xuất hiện mặt Ngửa. Tỉ số của số lần xuất hiện mặt Ngửa so với tổng số lần tung là
\choice
{$\dfrac{9}{20}$}
{\True $\dfrac{11}{20}$}
{$\dfrac{11}{9}$}
{$\dfrac{9}{11}$}
\loigiai{
Tỉ số là $\dfrac{11}{20}$.\\
\textbf{Đáp án của câu hỏi là B}
}
\end{ex}

\begin{ex} %[5X11N65]
Gieo một con xúc xắc $50$ lần, có $8$ lần xuất hiện mặt $6$ chấm. Tỉ số phần trăm số lần xuất hiện mặt $6$ chấm là
\choice
{\True $16\%$}
{$12\%$}
{$18\%$}
{$20\%$}
\loigiai{
Tỉ số phần trăm: $\dfrac{8}{50} = \dfrac{16}{100} = 16\%$.\\
\textbf{Đáp án của câu hỏi là A}
}
\end{ex}

\begin{ex} %[5X11H65]
Trong một trò chơi rút thẻ từ hộp, Nam rút $40$ lần và có $14$ lần rút được thẻ màu đỏ. Tỉ số của số lần rút được thẻ không phải màu đỏ so với tổng số lần rút là
\choice
{$\dfrac{7}{20}$}
{\True $\dfrac{13}{20}$}
{$\dfrac{7}{13}$}
{$\dfrac{13}{7}$}
\loigiai{
Số lần rút không phải màu đỏ: $40 - 14 = 26$ lần.\\
Tỉ số: $\dfrac{26}{40} = \dfrac{13}{20}$.\\
\textbf{Đáp án của câu hỏi là B}
}
\end{ex}

\begin{ex} %[5X11H65]
Bắn bia $25$ phát, có $20$ phát trúng tâm bia. Tỉ số phần trăm số phát bắn trúng tâm bia là
\choice
{$75\%$}
{$85\%$}
{$90\%$}
{\True $80\%$}
\loigiai{
$20 : 25 = 0{,}8 = 80\%$.\\
\textbf{Đáp án của câu hỏi là D}
}
\end{ex}

\begin{ex} %[5X11H65]
Quay vòng quay may mắn $30$ lần, mũi tên dừng lại ở ô “Chúc mừng” $6$ lần. Tỉ số của số lần dừng lại ở ô “Chúc mừng” so với tổng số lần quay viết dưới dạng phân số tối giản là
\choice
{\True $\dfrac{1}{5}$}
{$\dfrac{6}{30}$}
{$\dfrac{1}{6}$}
{$\dfrac{2}{5}$}
\loigiai{
$\dfrac{6}{30} = \dfrac{1}{5}$.\\
\textbf{Đáp án của câu hỏi là A}
}
\end{ex}

\begin{ex} %[5X11H65]
Mai thực hiện tung đồng xu $100$ lần thì thấy mặt Sấp xuất hiện $48$ lần. Tỉ số phần trăm số lần xuất hiện mặt Ngửa là
\choice
{$48\%$}
{$50\%$}
{\True $52\%$}
{$54\%$}
\loigiai{
Số lần mặt Ngửa: $100 - 48 = 52$ lần $\Rightarrow 52\%$.\\
\textbf{Đáp án của câu hỏi là C}
}
\end{ex}

\begin{ex} %[5X11V65]
Gieo một con xúc xắc $60$ lần, kết quả xuất hiện các mặt chẵn ($2, 4, 6$) là $33$ lần. Tỉ số của số lần xuất hiện mặt lẻ so với tổng số lần gieo là
\choice
{$\dfrac{11}{20}$}
{$\dfrac{33}{60}$}
{$\dfrac{1}{2}$}
{\True $\dfrac{9}{20}$}
\loigiai{
Số lần xuất hiện mặt lẻ: $60 - 33 = 27$ lần.\\
Tỉ số: $\dfrac{27}{60} = \dfrac{9}{20}$.\\
\textbf{Đáp án của câu hỏi là D}
}
\end{ex}

\begin{ex} %[5X11V65]
Một thủ môn bắt phạt đền $50$ lần, cản phá thành công $15$ lần. Tỉ số phần trăm số lần cản phá thành công của thủ môn đó là
\choice
{\True $30\%$}
{$25\%$}
{$35\%$}
{$40\%$}
\loigiai{
$15 : 50 = 30\%$.\\
\textbf{Đáp án của câu hỏi là A}
}
\end{ex}

\begin{ex} %[5X11V65]
Trong một hộp có $3$ quả bóng Xanh, Đỏ, Vàng. Rút ngẫu nhiên $1$ bóng rồi trả lại vào hộp. Sau $80$ lần rút, bóng Xanh xuất hiện $28$ lần, bóng Đỏ xuất hiện $32$ lần. Tỉ số phần trăm số lần rút được bóng Vàng là
\choice
{$20\%$}
{\True $25\%$}
{$30\%$}
{$35\%$}
\loigiai{
Số lần bóng Vàng: $80 - (28 + 32) = 20$ lần.\\
Tỉ số phần trăm: $20 : 80 = 0{,}25 = 25\%$.\\
\textbf{Đáp án của câu hỏi là B}
}
\end{ex}

\begin{ex} %[5X11V65]
Kiểm tra ngẫu nhiên $200$ sản phẩm của một dây chuyền may, phát hiện có $6$ sản phẩm lỗi. Tỉ số phần trăm sản phẩm đạt chất lượng tốt là
\choice
{$3\%$}
{$94\%$}
{\True $97\%$}
{$96\%$}
\loigiai{
Số sản phẩm tốt: $200 - 6 = 194$.\\
Tỉ số phần trăm: $194 : 200 = 97\%$.\\
\textbf{Đáp án của câu hỏi là C}
}
\end{ex}

\begin{ex} %[5X11V65]
Hai bạn An và Bình cùng chơi tung đồng xu. An tung $50$ lần được $24$ lần Sấp. Bình tung $50$ lần được $28$ lần Sấp. Tỉ số phần trăm số lần xuất hiện mặt Sấp của cả hai bạn gộp lại là
\choice
{$50\%$}
{$54\%$}
{$56\%$}
{\True $52\%$}
\loigiai{
Tổng số lần mặt Sấp: $24 + 28 = 52$ lần.\\
Tổng số lần tung: $50 + 50 = 100$ lần.\\
Tỉ số phần trăm: $52 : 100 = 52\%$.\\
\textbf{Đáp án của câu hỏi là D}
}
\end{ex}

\begin{ex} %[5X11C65]
Một xạ thủ bắn $100$ viên đạn vào bia. Tỉ số phần trăm số lần bắn trúng điểm 10 là $65\%$, bắn trúng điểm 9 là $25\%$, còn lại trúng điểm 8. Tổng số điểm mà xạ thủ đó đạt được sau 100 phát bắn là
\choice
{$945$ điểm}
{$960$ điểm}
{\True $955$ điểm}
{$950$ điểm}
\loigiai{
Số phát điểm 10: $65$ phát $\Rightarrow 65 \times 10 = 650$ điểm.\\
Số phát điểm 9: $25$ phát $\Rightarrow 25 \times 9 = 225$ điểm.\\
Số phát điểm 8: $100 - (65 + 25) = 10$ phát $\Rightarrow 10 \times 8 = 80$ điểm.\\
Tổng điểm: $650 + 225 + 80 = 955$ điểm.\\
\textbf{Đáp án của câu hỏi là C}
}
\end{ex}

\textbf{Phần 2. Trắc nghiệm đúng sai}
\setcounterpageref{ex}{0}

\begin{ex} %[5X11H65]
Thực hiện quay một tấm bìa hình tròn chia thành 4 phần bằng nhau ghi các số $1, 2, 3, 4$ tổng cộng $40$ lần. Kết quả dừng lại: Số 1: $8$ lần; Số 2: $12$ lần; Số 3: $14$ lần; Số 4: $6$ lần. Xét tính đúng sai:
\choiceTF
{\True Tỉ số số lần xuất hiện số 1 là $\dfrac{1}{5}$.}
{\True Tỉ số phần trăm số lần xuất hiện số 2 là $30\%$.}
{Số lần xuất hiện số lẻ chiếm $50\%$ tổng số lần quay.}
{\True Số lần xuất hiện số lẻ chiếm $55\%$ tổng số lần quay.}
\loigiai{
\begin{itemize}
    \item Số 1: $8 : 40 = \dfrac{1}{5}$. Đúng.
    \item Số 2: $12 : 40 = 30\%$. Đúng.
    \item Số lẻ (số 1 và số 3): $8 + 14 = 22$ lần $\Rightarrow 22 : 40 = 55\% \neq 50\%$. Ý 3 Sai, Ý 4 Đúng.
\end{itemize}
\begin{itemize}
    \item a) \textbf{Đúng}.
    \item b) \textbf{Đúng}.
    \item c) \textbf{Sai}.
    \item d) \textbf{Đúng}.
\end{itemize}
}
\end{ex}

\begin{ex} %[5X11V65]
Tung một con xúc xắc $120$ lần. Kết quả:
- Mặt 1 chấm: $18$ lần
- Mặt 2 chấm: $22$ lần
- Mặt 3 chấm: $20$ lần
- Mặt 4 chấm: $16$ lần
- Mặt 5 chấm: $24$ lần
- Mặt 6 chấm: $20$ lần
Xét tính đúng sai:
\choiceTF
{\True Tỉ số phần trăm mặt 5 chấm xuất hiện là $20\%$.}
{\True Tỉ số phần trăm mặt 1 chấm xuất hiện là $15\%$.}
{Tổng số lần xuất hiện mặt lớn hơn 4 chấm là $40$ lần.}
{\True Tổng số lần xuất hiện mặt lớn hơn 4 chấm là $44$ lần.}
\loigiai{
\begin{itemize}
    \item Mặt 5 chấm: $24 : 120 = 20\%$. Đúng.
    \item Mặt 1 chấm: $18 : 120 = 15\%$. Đúng.
    \item Mặt lớn hơn 4 chấm (5 và 6 chấm): $24 + 20 = 44$ lần $\neq 40$ lần. Ý 3 Sai, Ý 4 Đúng.
\end{itemize}
\begin{itemize}
    \item a) \textbf{Đúng}.
    \item b) \textbf{Đúng}.
    \item c) \textbf{Sai}.
    \item d) \textbf{Đúng}.
\end{itemize}
}
\end{ex}

\textbf{Phần 3. Trả lời ngắn}
\setcounterpageref{ex}{0}

\begin{ex} %[5X11H65]
Tung một đồng xu $50$ lần, có $27$ lần xuất hiện mặt Sấp. Tính tỉ số phần trăm số lần xuất hiện mặt Sấp. (Chỉ ghi số).
\shortans{54}
\loigiai{
\textbf{Đáp số:} $54$\\
$27 : 50 = 54\%$.
}
\end{ex}

\begin{ex} %[5X11H65]
Gieo một con xúc xắc $40$ lần, có $6$ lần xuất hiện mặt 6 chấm. Viết tỉ số của số lần xuất hiện mặt 6 chấm dưới dạng phân số tối giản $\dfrac{a}{b}$. Tính $a + b$.
\shortans{23}
\loigiai{
\textbf{Đáp số:} $23$\\
$\dfrac{6}{40} = \dfrac{3}{20} \Rightarrow a = 3, b = 20 \Rightarrow a + b = 23$.
}
\end{ex}

\begin{ex} %[5X11V65]
Một cầu thủ sút luân lưu $20$ lần, bóng vào lưới $16$ lần. Tính tỉ số phần trăm số lần sút bóng không vào lưới. (Chỉ ghi số).
\shortans{20}
\loigiai{
\textbf{Đáp số:} $20$\\
Không vào lưới: $4$ lần $\Rightarrow 4 : 20 = 20\%$.
}
\end{ex}

\begin{ex} %[5X11V65]
Bốc ngẫu nhiên một viên bi từ túi nhiều lần (có hoàn lại). Sau $60$ lần bốc, có $21$ lần bốc được bi xanh. Tính tỉ số phần trăm số lần bốc được bi xanh. (Chỉ ghi số).
\shortans{35}
\loigiai{
\textbf{Đáp số:} $35$\\
$21 : 60 = 0{,}35 = 35\%$.
}
\end{ex}

\begin{ex} %[5X11V65]
Một nhà máy kiểm tra $500$ sản phẩm thấy có $485$ sản phẩm đạt chuẩn. Tính tỉ số phần trăm sản phẩm không đạt chuẩn. (Dùng dấu phẩy nếu có).
\shortans{3}
\loigiai{
\textbf{Đáp số:} $3$\\
Không đạt chuẩn: $15$ sản phẩm $\Rightarrow 15 : 500 = 3\%$.
}
\end{ex}

\begin{ex} %[5X11C65]
Gieo một con xúc xắc $150$ lần, tỉ số phần trăm xuất hiện mặt số chẵn là $48\%$. Hỏi mặt số lẻ xuất hiện bao nhiêu lần?
\shortans{78}
\loigiai{
\textbf{Đáp số:} $78$\\
Mặt số lẻ chiếm $100\% - 48\% = 52\%$.\\
Số lần xuất hiện mặt lẻ: $150 \times 52 : 100 = 78$ lần.
}
\end{ex}
